% Document: TFG Avantprojecte - Object of the Project
% Language: English

\chapter{Introduction}

%\section{Problem Statement and Educational Gap}

Cybersecurity education faces a critical challenge: the disconnect between
theoretical knowledge of IT systems and practical understanding of how
security failures manifest in integrated environments. While students in
the Computer Engineering degree program learn fundamental concepts in databases,
network protocols, operating system internals, and software architecture,
they rarely experience how misconfigurations and design flaws in these domains
directly translate to security vulnerabilities.

Current training approaches amplify this gap. Commercial platforms such as
HackTheBox\cite{hackthebox} and TryHackMe\cite{tryhackme} focus on isolated challenge-based scenarios that
abstract away system complexity. Professional pentesting certifications
(e.g., CEH, OSCP\cite{ceh_framework,oscp_exam}) expect hands-on experience in realistic environments, yet
many academic programmes struggle to provide such environments due to cost,
complexity, or vendor lock-in.

Within the TecnoCampus Bachelor's Degree in Computer Engineering
(Management and Information Systems track), the ``Introduction to Cybersecurity''
course currently lacks a structured, reusable, and automatable lab environment that:

\begin{itemize}
    \item Aligns pentesting scenarios with specific degree learning outcomes
    \item Enables students to recognise vulnerabilities within integrated systems
          (not isolated exploits)
    \item Provides reproducible environments for consistent assessment
    \item Reduces instructor workload through automation and documentation
    \item Remains cost-effective and institution-independent (open-source,
          locally maintained)
\end{itemize}

\section{Context}

The Centre Universitari TecnoCampus~\cite{tecnocampus} offers applied science degrees
oriented towards professional practice in technology, business, creative industries,
and health. Affiliated with the Universitat Pompeu Fabra, it carries that
university's quality seal. Established in 2023 through the merger of three schools---the
Escola Superior Politecnica (ESUPT), the Escola Superior de Ciencies Socials i de
l'Empresa (ESCSET), and the Escola Superior de Ciencies de la Salut (ESCST)---it is
located at the technology park of the same name in Mataro, combining academic
infrastructure with an entrepreneurial ecosystem that facilitates student participation
in innovation and industry-facing projects.

The primary beneficiary of this project is the Bachelor's Degree in Computer
Engineering (Management and Information Systems track). Students in this programme
acquire foundational competencies in networks, operating systems, databases, and
software engineering all of which are directly relevant to security practice.
The \textit{Introduction to Cybersecurity} course is a natural integration point
for these skills, yet it currently lacks a dedicated hands-on environment.

Three perspectives illustrate the gap this project addresses:

\paragraph{Institutional perspective.}
The absence of dedicated cybersecurity lab infrastructure limits the institution's
capacity to offer competitive, hands-on security training aligned with degree
learning outcomes. Commercial alternatives such as HackTheBox and TryHackMe are
externally hosted, non-customisable, and not designed to map to specific academic
curricula or assessment rubrics.

\paragraph{Student perspective.}
Students encounter security concepts as isolated theoretical constructs rather than
as failures observable in realistic, integrated environments. Without exposure to
multi-component systems---firewalls, routers, servers, and applications working
together---they cannot develop the diagnostic reasoning required in professional
security roles.

\paragraph{Instructor perspective.}
Designing and maintaining practical security exercises is resource-intensive.
The absence of reusable, automatable lab environments forces instructors to
improvise or rely on third-party platforms, reducing pedagogical control and
limiting reproducibility across cohorts.

\section{Proposed Solution}

This project proposes a reusable penetration testing lab package built on
EVE-NG~\cite{eve_ng} and deployed on self-hosted infrastructure, designed
specifically for the \textit{Introduction to Cybersecurity} course at TecnoCampus.
The package comprises four progressive lab scenarios covering network reconnaissance
and enumeration, web application vulnerabilities, incident response and log forensics, and
privilege escalation. Labs~1--3 are fully implemented and validated; Lab~4
(privilege escalation) is in progress and will be completed in Phase~3. Each lab is self-contained: it includes an EVE-NG topology
file (\texttt{.unl}), a student guide (\textit{enunciado}) and an instructor
resolution guide (\textit{resolució}), an assessment rubric, and supporting
configuration and automation scripts. The entire package is open-source, vendor-
independent, and deployable on any institution running a compatible hypervisor.
Instructor deployment references, credentials, and lab-specific configuration
details are compiled in the Annexos (Volume~II), which accompanies this memory
and is also published online at \path{docs/web/docs/assets/official\_Documents/annexos.md}.
